\chapter{Grief Before Death}

\section*{Definition and Overview}
Grief before death, also called anticipatory grief, is the grief experienced by family members and loved ones while a person is still alive but dying, most often from a terminal illness. It involves mourning the losses that are already happening and the losses still to come: the person's declining health, the end of a shared future, and the approaching death. Anticipatory grief can include sadness, anger, guilt, anxiety, and exhaustion, and it can begin months before the death. It is a normal response, and understanding it can help carers and families cope, communicate, and make the most of the time they have together.

\section*{Epidemiology}
Anticipatory grief is experienced by a very large number of people. Most Australians die after a period of chronic or terminal illness, and their families face a period of caregiving and loss before death. Around half of Australians will provide informal care to a dying person at some stage, and carers report high rates of distress, anxiety, and depression during the terminal phase. Family caregivers commonly experience physical and emotional exhaustion, and anticipatory grief can coexist with, and sometimes complicate, the grief experienced after death.

\section*{Aetiology and Pathophysiology}
Anticipatory grief arises from the many losses that accompany a terminal illness.

\begin{itemize}
    \item \textbf{Ongoing losses:} the person's declining health and abilities, loss of independence, and loss of the relationship as it was
    \item \textbf{Anticipation:} knowing that death is coming, and the death of future hopes and plans
    \item \textbf{Caregiver strain:} the physical and emotional demands of caring can amplify grief
    \item \textbf{Role changes:} family members take on new roles, and roles and expectations change
    \item \textbf{Ambiguity:} uncertainty about the illness course can complicate grief
    \item \textbf{Individual differences:} people grieve the anticipated loss differently, depending on the relationship, culture, and supports
    \item \textbf{Both pre- and post-death:} anticipatory grief does not replace grief after death; it can coexist with it
\end{itemize}

\section*{Clinical Features}
Anticipatory grief presents with emotional, physical, and behavioural features.

\begin{itemize}
    \item Sadness and frequent crying
    \item Anxiety and worry about the future
    \item Anger, frustration, and irritability
    \item Guilt about not doing enough or feeling burdened
    \item Exhaustion and caregiver burnout
    \item Sleep disturbance and loss of appetite
    \item Difficulty concentrating and making decisions
    \item Withdrawal from others, or difficulty being with the dying person
    \item A sense of numbness or detachment
    \item Physical symptoms such as headaches and lowered immunity
\end{itemize}

\section*{Differential Diagnosis}
Anticipatory grief can overlap with conditions that warrant treatment.

\begin{itemize}
    \item Clinical depression, which is common in caregivers and should be treated
    \item Anxiety disorders and panic attacks
    \item Caregiver burnout, a state of physical and emotional exhaustion
    \item Post-traumatic stress responses to witnessing distressing illness
    \item Physical illness in the carer caused or worsened by stress
    \item Prolonged grief disorder developing after the death
\end{itemize}

\section*{When to Seek Medical Care}
Carers and family members should seek help if their own distress is affecting their health, ability to care, or daily function, or if they develop depression, anxiety, or thoughts of self-harm. Support services and counselling are available and beneficial.

\textbf{Call 000 (or local emergency services) for:}
\begin{itemize}
    \item Thoughts of harming yourself or ending your life
    \item Severe anxiety or panic that feels unmanageable
    \item A crisis in the care of the dying person, such as unmanageable pain, breathing difficulty, or sudden deterioration, which may need urgent palliative care input
    \item Any emergency health need for the carer or the person being cared for
\end{itemize}

\section*{Diagnosis and Investigations}
Anticipatory grief is recognised through sensitive conversation and assessment.

\begin{itemize}
    \item \textbf{Clinical interview:} exploring the family's experience, the illness course, and the supports in place
    \item \textbf{Assessment of the carer's wellbeing:} screening for depression, anxiety, exhaustion, and sleep problems
    \item \textbf{Assessment of supports:} practical, emotional, and spiritual supports available to the family
    \item \textbf{Review of the dying person's care:} ensuring palliative care needs are being met
    \item \textbf{Follow-up:} ongoing review of both the dying person and the family
\end{itemize}

\section*{Management}
Support for anticipatory grief focuses on the family's wellbeing while caring for the dying person.

\begin{itemize}
    \item \textbf{Palliative care:} good symptom control for the dying person reduces distress for everyone
    \item \textbf{Open communication:} encouraging honest conversations about the illness, wishes, and fears
    \item \textbf{Emotional support:} talking with family, friends, and support groups
    \item \textbf{Counselling:} grief and palliative care counsellors help families process anticipatory grief
    \item \textbf{Respite and practical help:} arranging respite care, home help, and shared caring duties
    \item \textbf{Self-care for carers:} protecting sleep, nutrition, and time away from caring
    \item \textbf{Completing unfinished business:} supporting meaningful time together, saying what needs to be said, and resolving conflicts
    \item \textbf{Planning:} advance care planning and funeral arrangements, which can reduce later burden
    \item \textbf{Medication:} treatment for depression or anxiety in the carer when present
\end{itemize}

\section*{Complications}
\begin{itemize}
    \item Caregiver burnout, with physical and emotional exhaustion
    \item Depression and anxiety in carers
    \item Complicated or prolonged grief after the death
    \item Family conflict over care decisions and roles
    \item Financial strain from time off work and care costs
    \item Physical health problems in carers from chronic stress
    \item Regret and guilt after the death if communication was difficult
\end{itemize}

\section*{Prognosis and Outcomes}
Anticipatory grief is normal, and most families navigate it with support. While it can be distressing, many people find that the period before death is also a time of closeness, meaning, and the opportunity to say goodbye. Families who communicate openly, use palliative care well, and support each other tend to cope better, both before and after the death. Support during the terminal phase also reduces the risk of complicated grief afterwards, and most carers find that grief resolves and life gradually reopens after the loss.

\section*{Prevention and Self-Care}
\begin{itemize}
    \item Acknowledge that anticipatory grief is normal and allow yourself to feel it
    \item Talk openly with the dying person about wishes, memories, and feelings
    \item Share caring duties with others rather than carrying the load alone
    \item Accept practical help with meals, cleaning, and transport
    \item Take regular breaks and protect time for yourself
    \item Maintain your own health: sleep, eat well, and see your GP if needed
    \item Use respite care and palliative support services
    \item Join a carers' support group
    \item Plan ahead with advance care planning and funeral arrangements when the time is right
    \item Seek counselling if grief is overwhelming or affecting your health
\end{itemize}

\section*{Sources and Further Reading}
The following reputable sources provide further information for healthcare students and patients seeking reliable, current advice about grief before death.

\begin{itemize}
    \item Palliative Care Australia. \textit{Support for carers and families.} https://palliativecare.org.au/
    \item CareSearch. \textit{Practical guidance for families facing terminal illness.} https://www.caresearch.com.au/
    \item GriefLine. \textit{Anticipatory grief support.} https://www.griefline.org.au/
    \item Carers Australia. \textit{Support and services for carers.} https://www.carersaustralia.com.au/
    \item Healthdirect Australia. \textit{Grief and loss.} https://www.healthdirect.gov.au/
    \item NHS. \textit{Anticipatory grief and supporting a dying relative.} https://www.nhs.uk/
\end{itemize}
